\appendix
\part{Приложения}

\chapter{Приложение А. Памятка по системам счисления и степеням двойки.}

\subsection*{Степени двойки}
\begin{center}
\begin{tabular}{c|c|c}
Степень & Значение & Где часто встречается \\
\hline
$2^8$ & $256$ & байт, диапазон \texttt{uint8\_t} \\
$2^{10}$ & $1024$ & килобайт в двоичной традиции \\
$2^{16}$ & $65536$ & диапазон 16-битного беззнакового числа \\
$2^{20}$ & $1048576$ & мегабайт в двоичной традиции \\
$2^{31}$ & $2147483648$ & граница 32-битного знакового \texttt{int} \\
$2^{32}$ & $4294967296$ & диапазон 32-битного беззнакового числа \\
$2^{64}$ & $18446744073709551616$ & диапазон 64-битного беззнакового числа \\
\end{tabular}
\end{center}

\subsection*{Быстрый перевод через группы битов}
\begin{itemize}
    \item $8 = 2^3$, поэтому одна восьмеричная цифра соответствует трём битам.
    \item $16 = 2^4$, поэтому одна шестнадцатеричная цифра соответствует четырём битам.
    \item Для перевода из двоичной системы в восьмеричную или шестнадцатеричную группы битов выделяются справа налево в целой части числа.
    \item Если в старшей группе не хватает битов, слева добавляются нули.
\end{itemize}

\chapter{Приложение Б. Краткий справочник команд Linux.}

\begin{center}
\begin{tabular}{l|p{9cm}}
Команда & Назначение \\
\hline
\texttt{pwd} & показать текущий каталог \\
\texttt{ls} & показать содержимое каталога \\
\texttt{cd} & перейти в другой каталог \\
\texttt{mkdir} & создать каталог \\
\texttt{touch} & создать пустой файл или обновить время изменения \\
\texttt{cp} & скопировать файл или каталог \\
\texttt{mv} & переместить или переименовать файл \\
\texttt{rm} & удалить файл \\
\texttt{less} & просмотреть текстовый файл постранично \\
\texttt{grep} & найти строки по шаблону \\
\texttt{find} & найти файлы в дереве каталогов \\
\texttt{chmod} & изменить права доступа \\
\texttt{man} & открыть справочную страницу \\
\end{tabular}
\end{center}

\subsection*{Стандартные потоки и перенаправления}
\begin{itemize}
    \item \texttt{stdin} -- стандартный ввод.
    \item \texttt{stdout} -- стандартный вывод.
    \item \texttt{stderr} -- стандартный поток ошибок.
    \item \texttt{command > file} -- записать стандартный вывод в файл.
    \item \texttt{command < file} -- подать содержимое файла на стандартный ввод.
    \item \texttt{command1 | command2} -- передать вывод первой команды на вход второй.
\end{itemize}

\chapter{Приложение В. Минимальные команды сборки C и NASM.}

\subsection*{C}
\begin{itemize}
    \item \texttt{gcc hello.c -o hello} -- собрать простую программу из одного файла.
    \item \texttt{gcc -c main.c -o main.o} -- получить объектный файл без линковки.
    \item \texttt{gcc main.o util.o -o program} -- слинковать несколько объектных файлов.
    \item \texttt{gcc -Wall -Wextra -g main.c -o main} -- собрать программу с предупреждениями и отладочной информацией.
\end{itemize}

\subsection*{NASM}
\begin{itemize}
    \item \texttt{nasm -f elf64 program.asm -o program.o} -- собрать 64-битный объектный файл для Linux.
    \item \texttt{ld program.o -o program} -- слинковать объектный файл напрямую через линковщик.
    \item \texttt{gcc program.o -o program} -- слинковать объектный файл через GCC, если программа использует стандартную библиотеку C.
\end{itemize}

\subsection*{Make}
\begin{itemize}
    \item \texttt{make} -- выполнить цель сборки по умолчанию.
    \item \texttt{make clean} -- выполнить цель очистки, если она описана в \texttt{Makefile}.
    \item \texttt{make target} -- выполнить конкретную цель.
\end{itemize}

\chapter{Приложение Г. Краткий справочник Git.}

\begin{center}
\begin{tabular}{l|p{9cm}}
Команда & Назначение \\
\hline
\texttt{git init} & создать локальный репозиторий \\
\texttt{git status} & посмотреть состояние рабочей директории \\
\texttt{git add file} & добавить файл в индекс \\
\texttt{git commit -m "message"} & создать коммит с сообщением \\
\texttt{git log} & посмотреть историю коммитов \\
\texttt{git diff} & посмотреть незакоммиченные изменения \\
\texttt{git clone url} & склонировать удалённый репозиторий \\
\texttt{git pull} & получить изменения из удалённого репозитория \\
\texttt{git push} & отправить локальные коммиты \\
\texttt{git branch} & посмотреть или создать ветки \\
\texttt{git switch branch} & переключиться на ветку \\
\texttt{git merge branch} & слить указанную ветку в текущую \\
\end{tabular}
\end{center}

\subsection*{Что обычно не нужно коммитить}
\begin{itemize}
    \item Исполняемые файлы, объектные файлы и временные результаты сборки.
    \item LaTeX-артефакты: \texttt{.aux}, \texttt{.log}, \texttt{.toc}, \texttt{.out}, \texttt{.fls}, \texttt{.fdb\_latexmk}.
    \item Локальные настройки редактора, если они не являются частью договорённостей проекта.
\end{itemize}
